\section{附录：Prompt 与 Tool 模板}

\subsection{Prompt Template 架构}

Burak 推荐的 prompt template 包含三个部分：

\begin{enumerate}
    \item Context：包含任务背景、policy 约束、用户意图
    \item Instruction：明确当前步骤的输出要求（比如「列出 top-3 解决方案」）
    \item Tool Usage：列出可选工具的调用方法，附带 success indicator
\end{enumerate}

\begin{knowledgebox}{Prompt Template 的检查清单}
每次更新 prompt 都要验证：
1. Context 是否更新（policy / data）
2. Instruction 是否仍然清晰
3. Tools \& examples 是否可复现
\end{knowledgebox}

\subsection{Tool Manifest 示例}

典型 tool manifest 包含以下字段：

\begin{itemize}
    \item tool name
    \item description
    \item input schema（参数类型、可选值）
    \item latency range
    \item safety constraints（哪些 user group 不能用）
\end{itemize}

即便 tool 只是内部 API，也建议写成 manifest，物流团队可以当作操作手册。

\subsection{Failure Tag 标签体系}

Burak 建议统一 failure tag taxonomy：

\begin{itemize}
    \item hallucination
    \item tool failure
    \item context expiration
    \item safety violation
    \item latency breach
\end{itemize}

每个 tag 关联处理流程：例如 `tool failure` 触发 fallback or escalate，`hallucination` 触发 LLM retrace path。

\subsection{本章小结}

Prompt 与 tool 模板是 agent reproducibility 的基石，把失败模式结构化之后，产品和工程可以用更小的迭代去修复。

